\documentclass[10pt,a4paper]{amsart}
\usepackage[margin=19mm]{geometry}
\usepackage{amsmath,amssymb,mathtools}
\usepackage[scr=rsfs,cal=euler]{mathalfa}
\usepackage{xfrac}
\usepackage[unicode,hidelinks,bookmarks=false]{hyperref}
\newcommand{\ie}{i.e.\ }
\newcommand{\cf}{cf.\ }
\newcommand{\resp}{respectively\ }
\newcommand{\Subsection}{\subsection}
\providecommand{\nobreakdash}{\nobreak\hbox{-}}
\setlength{\emergencystretch}{2em}
\multlinegap0pt
\usepackage{mathtools}
\usepackage{tikz}
\usepackage{xcolor}
\usepackage{amssymb}
\usepackage[shortlabels]{enumitem}
\newtheorem{proposition}{Proposition}
\def\S{{\mathbb S}}
\newcommand{\abs}[1]{\left |#1 \right |}
\newcommand{\aleq}{\precsim}
\newcommand{\brac}[1]{\left (#1 \right )}
\def\id{{\rm id\, }}
\def\n{{\mathcal N}}
\newcommand{\norm}[1]{\left\|{#1}\right\|}
\title{Generic non-uniqueness of minimizing harmonic maps from a ball to a sphere}
\author{Antoine Detaille, Katarzyna Mazowiecka}
\date{Source: arXiv:2403.12662v2 (CC BY 4.0)}
\begin{document}
\maketitle

\noindent\emph{Source and licence.} Generic non-uniqueness of minimizing harmonic maps from a ball to a sphere. Version arXiv:2403.12662v2 (CC BY 4.0), distributed by arXiv under the Creative Commons Attribution 4.0 International licence, http://creativecommons.org/licenses/by/4.0/, as recorded in the arXiv licence metadata for this exact version and shown on the abstract page https://arxiv.org/abs/2403.12662v2. Full licence notice: \texttt{licenses/arxiv-2403.12662-CC-BY-4.0.txt}.\par\smallskip

\noindent\textit{Editorial scope.} Counted unit: the article's proposition proposition:controlled\_homotopies (source0.tex lines 396--406). The article's complete original proof of it is delivered with it (source0.tex lines 424--509, one \texttt{proof} environment of 86 lines, which the article places away from the statement). No mathematics is written, repaired or re-derived: the statement and the proof are the article's own lines, in the article's own notation, and no retained line is substituted or re-flowed.  No further source-local context is needed: every cross-reference of every retained line resolves inside the two delivered spans.  Nothing was omitted from any retained span, so the package carries no omission record.  No retained line contains a citation, so the wrapper carries no bibliography and no bibliography entry.  No retained line contains a figure environment, an image-inclusion command or drawing code, so no image is delivered and the package declares no asset dependency.  Editorial formatting only, in addition: an AMS \texttt{amsart} preamble that restates the article's own declarations and macros used by the retained lines, namely the environments \texttt{proposition} on separate counters, together with the standard packages those lines need.  The retained environments print in this document as Proposition 1 in order of appearance, whereas the article prints the same items with the numbering of its own full document.  The complete native source of the article as distributed in the arXiv e-print for this exact version is bundled unchanged as \texttt{sources/156320/source0.tex}.
\begin{proposition}
\label{proposition:controlled_homotopies}
	Let $ \varphi \in C^{\infty}(\S^n,\n)$ and $p<n$.
	For every $ r > 0 $, for every $ x \in \S^n $, and every $ \psi \in C^{\infty}(\S^n,\n)$ homotopic to $ \varphi $ and satisfying $ \varphi = \psi $ on $ \S^n \setminus B_r(x)  $, there exists a homotopy $ H \in C^{\infty}(\S^n \times [0,1],\n)$ from $ \varphi $ to $ \psi $ such that
	\[
		\sup_{0 \leq t \leq 1} \norm{\varphi - H_{t}}_{W^{1,p}(\S^n)}
		\leq
		C\brac{\norm{\varphi}_{W^{1,p}(B_{2r}(x))} + \norm{\psi}_{W^{1,p}(B_{2r}(x))}}\text{,}
	\]
	for some constant $ C > 0 $ depending only on $ n $ and $ p $.
\end{proposition}

\begin{proof}[Proof of Proposition~\ref{proposition:controlled_homotopies}]
	Let $ G \in C^{\infty}(\S^n\times[0,1],\n)$ be any homotopy connecting $ \varphi $ to $ \psi $ with $G_0 = \varphi$ and $G_1 = \psi$.
	Since $ \varphi = \psi $ outside of $ B_r(x)  $, we may assume that $ G $ is stationary outside of $ B_r(x)  $, i.e., for each $t\in[0,1]$, we have $G_t= \varphi = \psi$ on $\S^n \setminus B_r(x)$.
	This claim can be proved with a by-hand construction, that we sketch below.
	We denote by \( \hat{x} \) the point at the antipode of \( x \).
	Let \( \Psi \colon \S^{n} \to \S^{n} \) be a smooth map such that \( \Psi = \id \) outside of \( B_{r}(x) \), and such that \( \Psi \) maps the annulus \( B_{r}(x) \setminus \overline{B}_{r/2}(x) \) diffeomorphically onto the annulus \( \S^{n} \setminus (\overline{B}_{r}(x) \cup \{\hat{x}\}) \), the circle \( \partial B_{r/2}(x) \) onto \( \{\hat{x}\} \), and the ball \( B_{r/2}(x) \) diffeomorphically onto \( \S^{n} \setminus \{\hat{x}\} \).
	It is readily observed that \( \Psi \sim \id \), through a homotopy stationary outside of \( B_{r}(x) \).
	Therefore, the maps \( u \circ \Psi \) and \( v \circ \Psi \) are homotopic to \( u \) and \( v \) respectively, through a homotopy stationary outside of \( B_{r}(x) \).
	Now, given a homotopy \( G' \) connecting \( u \) to \( v \), a homotopy \( G'' \) connecting \( u \circ \Psi \) to \( v \circ \Psi \) can be constructed by prescribing that \( G'' \) is stationary outside of \( B_{r} \), by letting \( G''_{t} = G'_{t} \circ \Psi \) on \( \overline{B}_{r/2} \) --- which corresponds to rescaling \( G' \) from \( \S^{n} \setminus \{ \hat{x} \}\) to \( B_{r}(x) \) --- and extending smoothly on the annulus \( B_{r} \setminus \overline{B}_{r/2} \).
	This is readily done by combining the observations that (i) \( u \circ \Psi \) and \( v \circ \Psi \) coincide also on \( B_{r} \setminus B_{r/2}(x) \) and are constant on \( \partial B_{r/2}(x) \), and (ii) \( G''_{t} \) is constant on \( \partial B_{r/2}(x) \).
	The required homotopy \( G \) stationary outside of $ B_r(x)  $ is then obtained by patching the three above homotopies, from \( u \) to \( u \circ \Psi \), from \( u \circ \Psi \) to \( v \circ \Psi \), and from \( v \circ \Psi \) to \( v \).
	
	Consider $ \tau > 0 $, which will be chosen sufficiently small at a later stage.
	We are going to rescale $ G $, $ \varphi $, and $ \psi $ from $ B_r(x)  $ to a smaller ball $ B_\tau(x) $, while keeping them unchanged outside of $ B_{2r}(x) $.
	More specifically, let $ \brac{\Phi_{t}}_{0 \leq t \leq 1} $ be a family of smooth diffeomorphisms of $ \S^n $ such that $ \Phi_{t} = \id $ outside of $ B_{2r}(x) $ and such that, on $ B_{2r}(x) $, in the local chart given by the exponential map around $ x $, $ \Phi_{t} $ is expressed as
	\[
		\begin{cases}
			\frac{rx}{\brac{1-t}r+t\tau} & \text{if $ \abs{x} \leq \brac{1-t}r+t\tau $,} \\
			\frac{x}{\abs{x}}\brac{\frac{r}{2r-\brac{1-t}r-t\tau}\brac{\abs{x}-\brac{1-t}r-t\tau}+r} & \text{if $ \brac{1-t}r+t\tau \leq \abs{x} \leq 2r $.}
		\end{cases}
	\]
	We define $ H \in C^{\infty}(\S^n\times[0,1],\n)$ by
	\[
		H_{t} \coloneqq
		\begin{cases}
			\varphi \circ \Phi_{3t} & \text{if $ 0 \leq t \leq \frac{1}{3} $,} \\
			G_{3\brac{t-1/3}} \circ \Phi_{1} & \text{if $ \frac{1}{3} \leq t \leq \frac{2}{3} $,} \\
			\psi \circ \Phi_{1-3\brac{t-2/3}} & \text{if $ \frac{2}{3} \leq t \leq 1$.}
		\end{cases}
	\]
	Of course, $ H $ is a homotopy from $ \varphi $ to $ \psi $.
	It remains to show that, if $ \tau > 0 $ is suitably small, then $ H $ satisfies the required estimate.

	For $ 0 \leq t \leq \frac{1}{3} $, we note that $\varphi - H_t =0$ outside $B_{2r}(x)$. We readily obtain bounds on the Jacobian and the derivatives of $ \Phi_{t} $, so that the change of variable theorem combined with $n-p>0$ implies that
	\[
			\norm{\varphi - H_{t}}_{W^{1,p}(\S^n)}
		\leq
		\norm{\varphi}_{W^{1,p}(B_{2r}(x))} + \norm{\varphi \circ \Phi_{3t}}_{W^{1,p}(B_{2r}(x))}
		\aleq \norm{\varphi}_{W^{1,p}(B_{2r}(x))}\text{.}
	\]
	Similarly, for $ \frac{2}{3} \leq t \leq 1 $, we have
	\[
		\norm{\varphi - H_{t}}_{W^{1,p}(\S^n)}
		\leq
		\norm{\varphi}_{W^{1,p}(B_{2r}(x))} + \norm{\psi \circ \Phi_{3t}}_{W^{1,p}(B_{2r}(x))}
		\lesssim
		\norm{\varphi}_{W^{1,p}(B_{2r}(x))} + \norm{\psi}_{W^{1,p}(B_{2r}(x))}\text{.}
	\]
	Concerning $ \frac{1}{3} \leq t \leq \frac{2}{3} $, we estimate
	\[
	\begin{split}
		\norm{\varphi - H_{t}}_{W^{1,p}(\S^n)}
		&\leq
 		\norm{\varphi}_{W^{1,p}(B_{2r}(x))} + \norm{G_{3\brac{t-1/3}} \circ \Phi_{1}}_{W^{1,p}(B_{2r}(x))}\\
 		&\lesssim
  		\norm{\varphi}_{W^{1,p}(B_{2r}(x))} + \norm{G_{3\brac{t-1/3}}}_{W^{1,p}(B_{2r}(x) \setminus B_r(x)) } + \tau^{\frac{n-p}{p}}\norm{G_{3\brac{t-1/3}}}_{W^{1,p}(B_{2r}(x))}\text{.}
	\end{split}
	\]
	Since the homotopy $ G $ has been assumed to be stationary outside of $ B_r(x)  $, we know that $ \norm{G_{3\brac{t-1/3}}}_{W^{1,p}(B_{2r}(x) \setminus B_r(x))} = \norm{\varphi}_{W^{1,p}(B_{2r}(x)\setminus B_r(x) )} $.
	On the other hand, by compactness, we have
	\[
		\sup_{0 \leq t \leq 1} \norm{G_{t}}_{W^{1,p}(B_{2r}(x))}
		\leq
		C_1
	\]
	for some possibly large constant $ C_1 > 0 $.
	We may assume that either \( \norm{\varphi}_{W^{1,p}(B_{2r}(x))} \neq 0 \) or \( \norm{\psi}_{W^{1,p}(B_{2r}(x))} \neq 0 \).
	Indeed, if \( \norm{\varphi}_{W^{1,p}(B_{2r}(x))} = 0 = \norm{\psi}_{W^{1,p}(B_{2r}(x))} \), this implies that both \( \varphi \) and \( \psi \) are identically zero --- note that this may only happen if \( 0 \in \n \) --- and we may directly conclude by choosing \( H \) to be constantly zero.
	As $ p < n $, we may therefore choose $ \tau > 0 $ sufficiently small, depending on $ C_1 $, so that
	\[
		\tau^{\frac{n-p}{p}}\norm{G_{3\brac{t-1/3}}}_{W^{1,p}(B_{2r}(x))}
		\leq
		\norm{\varphi}_{W^{1,p}(B_{2r}(x))} + \norm{\psi}_{W^{1,p}(B_{2r}(x))}
		\quad
		\text{for every $ \frac{1}{3} \leq t \leq \frac{2}{3} $.}
	\]
	Hence, we deduce that
	\[
		\norm{\varphi - H_{t}}_{W^{1,p}(\S^n)}
		\lesssim
		\norm{\varphi}_{W^{1,p}(B_{2r}(x))} + \norm{\psi}_{W^{1,p}(B_{2r}(x))}
		\quad
		\text{for every $ \frac{1}{3} \leq t \leq \frac{2}{3} $.}
	\]
	This concludes the proof.
\end{proof}

\end{document}
